\chapter{The Finite Gauge}
% OUTLINE Ch2 "The Finite Gauge". Two moves kept distinct from Ch1 (Kodo's
% constraint): Ch1 READ the number; Ch2 is the GAUGE that reads it + why
% FINITE is earned. Content: the verify-API (verify(claimed)->{verified?,
% residual}) as a gauge; the four books = four gauges; the reading as a RATIO
% of counts, no scale trusted (Vol 1 Wheatstone/null); FINITE = bounded by
% construction, the total fragment that always halts (Vol 2 s7.3 Turing-point,
% count-to-3-and-stop). "gauge" = the METROLOGICAL sense (a measuring
% instrument), NOT gauge-theory/QED (that door stays shut). Bold-math/
% fenced-referent continuous; no Lean identifiers; step-not-move. Gate: Kodo.

The last chapter read a number; this one says what kind of instrument
reads a number that way, because the title of this volume makes two claims
about the instrument --- that it is a \emph{gauge}, and that it is
\emph{finite} --- and both are earned here or not at all.

Take \emph{gauge} first, in the plain metrological sense: a gauge is a
tool that reads a quantity by comparison, not by trusting an absolute
scale. A pressure gauge reads against a spring; a bridge, the first volume's
diamond, reads a resistance against three others and believes only the
balance. This series' instrument is a gauge in exactly that sense, and the
form it takes is a single act named at the instrument's own specification
level: hand the instrument a \emph{claimed} value and it returns not a
verdict from on high but a \emph{residual} --- the signed distance between
what was claimed and what the model's own counts actually read. Verify a
claim; receive a difference. That is the whole interface, and it is a gauge
to the bone: it never announces the model's number as an oracle would; it
holds up a claim against the model's counts and reports the gap. The number
the last chapter stated is what the residual measures against, and it is
the model's own throughout --- the gauge reads the model's coupling, never
the world's, and returns a residual in the model's own units.

The series has four such gauges, one per volume, reading the same model
through four descriptions --- geometry, intuition, grammar, and the
measurement that grounds them --- and this volume is where their common
reading is stated. What makes them one instrument rather than four is that
each returns its residual in counts of the same kind, comparable because
they are all the model's own, none imported from a scale outside. A reading
that trusts no external standard cannot drift against one; that is the
gauge's honesty, and it is the reason the number of the last chapter is an
exact ratio and not a rounded measurement. The instrument reads a ratio of
its own counts and believes only the comparison --- the bridge's
discipline, four volumes on.

Now \emph{finite}, which is the word this volume's title stakes the most
on, and it is earned by a property the second volume of this series proved
about the instrument and named there exactly. The gauge is not merely
small; it is \emph{total} --- every reading it performs halts, on every
lawful input, by construction. There is no claim it can be handed on which
it computes forever; the machinery that would let a program run without end
is deliberately absent from every part of the instrument that carries a
reading, quarantined where nothing load-bearing can reach it. So a residual
is always delivered: ask the gauge to verify a claim and it finishes and
hands back the difference, always, in bounded time. That is what
\emph{finite} means in the title --- not that the answer is a small number,
but that the reading is a terminating computation, a walk that always
reaches its stop. The single invariant is bracketed in a \emph{finite}
number of earned steps --- three runs and a halt --- because the gauge that
reads it cannot do otherwise.

And here the two words in the title meet, because finiteness is what makes
the gauge honest about what it does not know. An instrument that could run
forever could always promise that one more step would close the gap to any
target you named; the promise would be empty, but it would be available, and
availability is exactly what a number under suspicion does not need. A finite
gauge cannot make that promise. It runs its earned steps, it halts, and it
reports the bracket it actually reached --- never a point it might reach if
allowed to run on. The bracket-not-point discipline the whole series keeps
is not a stylistic vow; it is what a terminating instrument is physically
able to say. The gauge stops because it must, and because it stops, the most
it can honestly report is walls, not a landing. Finiteness and honesty are
the same property of the machine, read from two sides.

One boundary closes the chapter, and it guards the title's word from a
reading it does not carry. \emph{Gauge} here is the metrologist's word for a
comparing instrument, and nothing more; it is not the physicist's
gauge --- the freedom-and-symmetry apparatus of field theory --- and this
volume makes no claim in that register. The instrument compares the model's
counts against a claim and returns a difference; it does not gauge a field,
fix a symmetry, or touch the machinery by which the world's forces are
written. The word is borrowed from the workshop, not the field theory, and
it is returned unspent in the other sense. What the finite gauge reads is
the model's single invariant, and how far that reading rhymes with problems
far larger than itself --- while touching none of them --- is the next
chapter's careful business.
